\section{Introduction}
\label{sec:intro}

There is a quiet asymmetry in applied financial econometrics. Producing a result
is cheap; making that result re-runnable, by someone else, on demand, is
expensive, and so it is mostly not done. A reader of a typical spillover or
contagion paper is handed a table and an appendix, not a button. When the
printed coefficient is later questioned, the burden of reconstruction falls on
whoever raises the question, and the original computation, with its package
versions, its random seeds, and its data vintage, is rarely recoverable in full.

We take the view that this is, for the most part, a software problem, and that
it is worth attacking with the engineering discipline that a serious systems
group would bring to any service whose correctness matters. Three failure modes
recur. The first is undocumented data revisions: a vendor restates a close,
a roll convention shifts a futures series, and a number that was correct last
quarter no longer reproduces, with nothing in the published record to say so.
The second is estimator detail buried in analysis scripts: an embedding lag, a
neighbour count, a bandwidth, a partition of the sample into episodes, choices
that move a headline figure but live only in a working directory. The third is
machine-specific numerics: an unpinned multi-threaded linear-algebra library, a
nearest-neighbour tie broken differently on a different processor, a sum
accumulated in a different order, each capable of moving the third decimal place
in a way no reader can anticipate. None of these is exotic. Together they are
why a result in this field is so often easy to publish and hard to re-run.

\subsection{Thesis: one engine, one substrate, one gate}

This manuscript describes our response, which is deliberately a programme and not
a paper. \engine{} is organised as a research ``space programme'': one sandboxed
compute engine, exposing one methodological substrate of eight primitives, under
one pre-registered and failable evaluation gate. The six published frameworks
that the group has produced, on adaptive market networks (NAMH), scale-ordered
contagion (SOCH), cross-border contagion channels, wavelet quantile transfer
(WaveQTE), commodity systems, and protocol-mediated systemic risk (MCPFM), are
\emph{elements} of that one programme, not separate codebases that happen to
share an author. The same long-memory estimator, the same wavelet decomposition,
the same transfer-entropy kernel, the same surrogate machinery, and the same
community detection thread through all of them. Verifying the substrate, once,
carefully, verifies every element built on it.

The engine itself is small. It runs a fixed catalogue of econometric methods
inside a sandbox and returns each result as structured data, with enough
provenance attached that the number can be banded against an expectation and
reproduced later. As of the live release it serves twenty-six methods, and the
most recent full suite recorded twenty-six of twenty-six passing. It is written
in Node and R, deployed for the research platform of the School of Humanities,
Social Sciences and Management at IIT Bhubaneswar, and documented from the engine
outward in the published reference \cite{bhandari2026econstellar}. The design
thesis is one sentence: \emph{a method is not in the catalogue until it ships
with a failable test of its own output.} The catalogue and that suite of tests
are the deliverable; the elements are what the deliverable certifies.

\subsection{Engine, not library}

It would be natural to read what follows as a library of econometric routines.
It is not, and the distinction is load-bearing both for safety and for trust.

A library accepts arbitrary code and runs it; the surface for arbitrary-code
execution is the point. The engine accepts no code. A request names a method
from a fixed registry and a parameter object validated against a schema, and the
runner scripts that do the work are part of the repository, never user-supplied.
\textbf{The parameterised-only registry is therefore the primary guard against
arbitrary-code execution}, and the operating-system sandbox is defence in depth
on top of it, not the load-bearing control. On the workstation and in local
development each analysis runs under \code{bubblewrap} with the network
unshared, the filesystem read-only, and a fresh scratch space; the deployed
cloud image ships without that tool and runs under a wall-clock timeout with the
managed runtime as its kernel boundary. Either way, no user input reaches a
shell, because the registry admits no user code in the first place. We return to
this topology, and to the three planes the engine runs across, in
Section~\ref{sec:architecture}.

The other half of ``engine, not library'' is the catalogue rule already stated:
a method joins only with a pre-registered, failable test of its own output, and
holes in coverage are marked rather than filled. The registered voice of the
project is blunt about the trade this implies. \emph{We would rather have
twenty-six checked methods than fifty asserted ones.} A figure that the engine
reports is a figure the engine can fail to reproduce, loudly and in public, and
that is the property we are trying to buy.

\subsection{Discipline borrowed from systems practice, applied to econometrics}

The engineering posture here is borrowed, openly, from how top-tier systems and
research-infrastructure groups treat reliability and verification: correctness
asserted by a pre-registered evaluation rather than by the authors, every figure
resolving to a job identifier or a revision string, machine-written ledgers that
are never hand-edited, and honest reporting of what does not work. What is
unusual is the domain. We apply that discipline to financial econometrics, a
field whose objects are noisy, whose data is revised, and whose estimators are
sensitive to choices that rarely survive into print. The combination is the
contribution: not any single number, but the apparatus that makes each number
re-runnable and the negatives as visible as the positives.

That apparatus is needed precisely because of what the markets are. Under the
adaptive markets hypothesis \citep{lo2004}, efficiency is not a fixed property
but an evolutionary one: heterogeneous participants adapt at different rates,
relationships strengthen and decay, and the structure a study measures in one
regime need not hold in the next. The Grossman--Stiglitz argument
\citep{grossman1980} sharpens the point from the side of information: if prices
revealed all private information costlessly, no one would pay to gather it, so
informationally efficient markets are impossible and a residue of private,
costly, unequally held information must remain. Markets that are adaptive,
heterogeneous, and information-driven in this sense do not yield to a single
estimate stated once. They demand a measurement apparatus that can be re-run as
the regime turns, that states its full parameterisation, and that reports
honestly when a structure it found before has dissolved. The engine is our
attempt to build exactly that apparatus and to hold it to a standard a reader can
check.

\subsection{Contributions}

This manuscript makes the following contributions.

\begin{enumerate}[leftmargin=1.6em]
  \item \textbf{A programme framing.} We argue that six separately published
  frameworks are better understood, and better maintained, as elements of one
  research programme sharing a single eight-primitive substrate and a single
  evaluation gate, and we make that claim concrete by exhibiting which primitives
  each element instantiates (Section~\ref{sec:elements}).

  \item \textbf{An engine whose safety rests on parameterisation.} We describe a
  compute engine in which the parameterised-only registry, not the sandbox, is
  the primary control against arbitrary-code execution, with the sandbox as
  defence in depth (Section~\ref{sec:architecture}).

  \item \textbf{A catalogue under a failable gate.} We state and apply the rule
  that a method enters the catalogue only with a pre-registered, failable test of
  its own output, served live as twenty-six methods at twenty-six of twenty-six
  passing, and we describe the evaluation discipline that distinguishes a band
  failure from a harness failure (Section~\ref{sec:elements} and the
  verification section).

  \item \textbf{Governed determinism and bit-exact acceleration.} We show that a
  governed core budget with pinned linear-algebra threading makes a result
  independent of worker count, and that a GPU port of the transfer-entropy
  nearest-neighbour search reproduces the CPU estimate to within machine epsilon
  while running several times faster (the heterogeneous-compute section).

  \item \textbf{Honest negatives as first-class results.} We report the negatives
  alongside the positives, including a contagion significance test that is not
  significant, a false-discovery-controlled network that is degenerate, and a
  pre-registered model re-evaluation whose discrimination falls below chance, and
  we treat each as an immutable fact of the record (Sections~\ref{sec:elements}
  and~\ref{sec:outlook}).
\end{enumerate}

\subsection{Related work}

Three literatures meet in this engine. The first is the long concern with
reproducibility in empirical economics, opened by the \emph{Journal of Money,
Credit and Banking} replication project~\cite{dewald1986replication} and sharpened
by studies of the numerical reliability of econometric software~\cite{mccullough1999numerical}
and of reproducible econometric practice~\cite{koenker2009reproducible}. The
computational-science community has codified much of the operational discipline we
adopt: simple rules for reproducible research~\cite{sandve2013ten}, container-based
environments~\cite{boettiger2015docker}, and the older idea of literate,
provenance-carrying programs~\cite{knuth1984literate}. Our addition to this lineage
is not another guideline but a running engine in which the discipline is enforced
by a failable gate.

The second is the information-theoretic and connectedness toolkit the substrate
implements. Transfer entropy~\cite{schreiber2000} and its nearest-neighbour
estimation~\cite{kraskov2004,frenzel2007} underpin the engine's directional-flow
methods; effective transfer entropy and surrogate correction~\cite{marschinski2002,schreiber2000surrogate}
control bias, applied financial studies~\cite{dimpfl2013using} motivate the design,
and detrended fluctuation analysis~\cite{peng1994mosaic} is the long-memory
primitive. The connectedness family follows the Diebold--Yilmaz
programme~\cite{diebold2009measuring,diebold2012,diebold2014} and its frequency
extension~\cite{barunik2018}, and we read the systemic-risk feed alongside the
standard market measures: CoVaR~\cite{adrian2016covar}, systemic expected
shortfall~\cite{acharya2017measuring}, and SRISK~\cite{brownlees2017srisk}.

The third is the view of market efficiency the elements speak to: the
efficient-markets benchmark~\cite{fama1970efficient}, the impossibility of
informationally efficient markets~\cite{grossman1980}, heterogeneous-agent
dynamics~\cite{brock1998heterogeneous}, and the adaptive-markets
synthesis~\cite{lo2004} that gives the NAMH element its name.

\subsection{Roadmap}

The paper proceeds from the engine outward. Section~\ref{sec:architecture} sets out the
three-plane architecture, the request lifecycle, and the sandbox topology. The
substrate section catalogues the eight primitives and the twenty-six methods
built on them, and the heterogeneous-compute and AI-layer sections cover the
governed core budget, the GPU offload held to a bit-exact bar, and the
natural-language front door onto the same validated registry. The verification
and epistemics section states how a pre-registered, failable evaluation suite,
a machine-written state ledger, and a claims record that never deletes combine
into a self-checking system, with worked examples following.
Section~\ref{sec:elements}, \emph{The Elements}, then reads the six published
frameworks as elements of the one programme, each with its economic question,
the substrate primitives it instantiates, its headline verified result, and its
honest status. Section~\ref{sec:outlook} distils the design principles, states
the limitations plainly, and describes the self-renewing programme we are
building towards.
